\chapter{Conclusão}\label{cap:conclusao}

Este trabalho apresentou uma investigação experimental sobre o diagnóstico de falhas incipientes em transformadores de potência por Análise de Gases Dissolvidos (DGA). O estudo combinou redes neurais LSTM para previsão temporal de concentrações gasosas, árvores com aumento de gradiente (XGBoost) para classificação dos estados de integridade, busca evolutiva por meio do Algoritmo Genético Melhorado (IGA) para seleção de variáveis e ajuste de hiperparâmetros, e análise de interpretabilidade com TreeSHAP. A pesquisa buscou avaliar empiricamente a real eficácia de arquiteturas híbridas e o impacto da inclusão de previsões temporais frente aos métodos normativos tradicionais e às características dos dados práticos de campo.

\section{Síntese das Contribuições}\label{sec:sintese_contribuicoes}

A primeira contribuição metodológica do trabalho foi a estruturação da análise em dois eixos independentes de investigação (Track A e Track B), respeitando as naturezas distintas das bases de dados disponíveis na literatura de DGA. A verificação prévia e a deduplicação dos dados removeram cerca de 15\% de amostras repetidas na base estática, assegurando que as partições de treino e teste fossem estritamente independentes e evitando estimativas infladas de desempenho.

No aspecto da otimização algorítmica, o acoplamento do IGA ao classificador XGBoost mostrou que o principal benefício da abordagem evolutiva esteve na simplificação do modelo. No Track A, o IGA reduziu o conjunto de variáveis de 17 para 8 atributos (uma redução de 53\%), eliminando entradas colineares e mantendo um escore $F_1$-macro idêntico de 0,891. Essa redução no número de variáveis é relevante para a implementação prática em sistemas de monitoramento contínuo em subestações, pois diminui o esforço computacional e reduz a dependência de múltiplos canais de sensoriamento.

No diagnóstico multiclasse do Track B (sete classes da norma IEC 60599), a avaliação por validação cruzada repetida em 25 dobras revelou um patamar estável de desempenho entre $F_1 \approx 0,75$ e $0,76$. Os modelos baseados em XGBoost (tanto o proposto quanto o de referência da literatura) superaram amplamente os métodos normativos rígidos (Triângulo 1 de Duval com $F_1 = 0,363$ e Razões IEC 60599 com $F_1 = 0,257$). No entanto, os diferentes algoritmos de aprendizado supervisionado apresentaram desempenhos estatisticamente equivalentes entre si, apontando para um limite prático imposto pelo uso exclusivo dos cinco gases convencionais e pela sobreposição física existente entre tipos vizinhos de falha.

Por fim, a análise de interpretabilidade com o algoritmo TreeSHAP permitiu verificar a coerência físico-química das decisões aprendidas pelo XGBoost. O modelo identificou de forma consistente o acetileno como o principal indicador de descargas elétricas e o etileno como marcador de falhas térmicas, alinhando-se aos princípios da termodinâmica de degradação do óleo. Além disso, a extração de limiares marginais identificou que a transição para descargas de baixa energia ocorre na faixa entre 20\% e 25\% de etileno, o que coincide de maneira direta com o limite de 23\% estabelecido empiricamente no Triângulo 1 de Duval.

\section{Avaliação da Hipótese Temporal e Resultados Negativos}\label{sec:defesa_resultados_negativos}

Um resultado relevante desta pesquisa foi a constatação empírica de que o fornecimento de concentrações futuras estimadas pela rede LSTM não gerou ganho estatisticamente significativo no diagnóstico pontual de falhas no Track A. O contraste pareado por \textit{bootstrap} com 5.000 iterações indicou uma variação marginal de apenas $\Delta F_1 = +0,0056$, com intervalo de confiança de 95\% entre $-0,0037$ e $+0,0150$ ($p \approx 0,238$). Como o intervalo engloba o zero ($p > 0,05$), os modelos com e sem projeções da LSTM mostraram-se estatisticamente equivalentes.

A razão física e matemática para essa equivalência reside no comportamento de longo prazo dos transformadores analisados. Em equipamentos sob regime operacional contínuo, a evolução das concentrações ocorre de forma gradual, de modo que as leituras atuais e as projeções de curto prazo apresentam coeficientes de correlação superiores a 0,997. Consequentemente, as previsões futuras fornecidas pela rede neural atuam na prática como variáveis fortemente redundantes em relação às medições observadas no momento, comportamento confirmado pela migração de crédito observada nos valores de Shapley.

O resultado negativo documentado mostra que, para a base avaliada, o acoplamento da rede recorrente não justifica o custo computacional adicional para o diagnóstico pontual. Enquanto a LSTM modela a evolução quantitativa das séries ($R^2 > 0,99$ em curto prazo), as projeções atuam como variáveis fortemente colineares às leituras atuais ($r > 0,997$). Dessa forma, o classificador alimentado unicamente pelas medições do momento atinge desempenho equivalente com arquitetura mais simples e menor esforço de implantação.

\section{Limitações da Pesquisa}\label{sec:limitacoes_pesquisa}

Os resultados e as conclusões apresentadas devem ser interpretados considerando as limitações dos dados públicos disponíveis na literatura de DGA:

\begin{enumerate}
    \item \textbf{Fragmentação das bases de dados:} Não se identificou na literatura aberta um conjunto de dados que una simultaneamente séries temporais contínuas em alta taxa de amostragem, medições completas dos sete gases normativos e laudos confirmados de defeitos severos. Essa indisponibilidade motivou a divisão do estudo nos dois eixos investigados.
    \item \textbf{Escopo restrito no Track A:} O conjunto temporal de \citeonline{katser2022} disponibiliza apenas quatro gases ($\text{H}_2$, $\text{CO}$, $\text{C}_2\text{H}_4$ e $\text{C}_2\text{H}_2$), impossibilitando o cálculo de razões clássicas que envolvem metano e etano. Adicionalmente, as concentrações registradas são predominantemente baixas (na escala de $10^{-3}$), refletindo condições de usinas nucleares com monitoramento rigoroso e pouca ocorrência de falhas destrutivas de grande porte.
    \item \textbf{Heterogeneidade e ruído de rotulagem no Track B:} A base estática compilada reúne ensaios laboratoriais de diferentes concessionárias e períodos. Conforme demonstrado na auditoria normativa, parte das amostras rotuladas como defeitos térmicos severos (T3) apresenta características cromatográficas incompatíveis com os critérios normativos internacionais, gerando um ruído de anotação que limita o desempenho máximo que qualquer modelo supervisionado pode alcançar sobre esses dados.
\end{enumerate}

\section{Direcionamentos para Trabalhos Futuros}\label{sec:trabalhos_futuros}

A partir das constatações e limitações identificadas neste estudo, sugerem-se os seguintes caminhos para continuidade da pesquisa:

\begin{enumerate}
    \item \textbf{Estimativa de tempo até a falha e vida útil remanescente:} Em vez de utilizar as projeções da rede LSTM para alimentar um classificador de estado pontual, a capacidade preditiva temporal pode ser direcionada para estimar quando as concentrações de gases cruzarão os limites críticos de atenção das normas técnicas, auxiliando diretamente o planejamento de paradas de manutenção.
    \item \textbf{Análise da taxa de variação dos gases:} Desenvolver modelos focados na velocidade diária ou mensal de geração dos gases ($\Delta\text{ppm}/\text{tempo}$), que na prática das concessionárias é um indicador de severidade frequentemente mais sensível do que as concentrações absolutas isoladas.
    \item \textbf{Integração com variáveis operacionais externas:} Incorporar variáveis de operação dos transformadores, como histórico de carregamento elétrico, temperatura do topo do óleo e umidade, permitindo contextualizar a produção de gases com o estresse térmico e elétrico suportado pelo equipamento.
    \item \textbf{Validação com bases temporais completas:} Conduzir novos testes caso venham a ser disponibilizados conjuntos de dados temporais de campo que incluam o espectro completo dos gases combustíveis e registros de equipamentos submetidos a diferentes regimes industriais.
\end{enumerate}
